% Generated by scripts/blueprint_gen.py; do not edit.
\chapter{Proof of H1: the Dolbeault–Grothendieck lemma}\label{chap:Z4}
The plan proves the lemma first on open subsets of ℂⁿ by Grothendieck's induction, then transfers it to X through holomorphic charts. The analytic core reduces to one parametric Cauchy transform, built from TauCeti's Cauchy–Pompeiu file and Mathlib's parametric convolution lemmas.

\textbf{Target.} For every open U ⊆ X, q ∈ ℕ, x ∈ U and ω ∈ Γ(U, 𝒜\textsuperscript{0,q+1}) with ∂̄ω = 0, there are an open U′ with x ∈ U′ ⊆ U and η ∈ Γ(U′, 𝒜\textsuperscript{0,q}) with ∂̄η = ω|U′.

\textbf{Four layers.}

\begin{enumerate}
\item \emph{Parametric ∂̄-solver in one variable} (pure analysis, Banach-valued). Solves ∂g/∂z̄\textsubscript{k} = f near x, keeps g holomorphic in every variable in which f is, and commutes with continuous linear maps.
\item \emph{Flat Dolbeault complex on ℂⁿ}. Forms are maps into Mathlib's \texttt{ℂⁿ [⋀\textasciicircum{}Fin q]→L[ℝ] ℂ}; ∂̄ω = Σ\textsubscript{j} dz̄\textsubscript{j} ∧ ∂̄\textsubscript{j} ω. Contains the induction and the flat lemma.
\item \emph{Comparison with TauCeti}. Identifies TauCeti's \texttt{delBar} on open subsets of ℂⁿ with the flat ∂̄.
\item \emph{Transfer to X} through charts and \texttt{mpullback\_delBar}, then instantiates \texttt{DolbeaultGrothendieck E X}.
\end{enumerate}
\textbf{Two simplifications drive the design.} The statement is local, so each induction step may shrink the neighbourhood; Hörmander's exhaustion argument for a whole polydisc is not needed. The sup norm on \texttt{Fin n → ℂ} makes \texttt{Metric.ball x r} a polydisc, so no polydisc API is needed.

\textbf{Upstream route.} TauCeti's roadmap ComplexKahlerGeometry, stage K2, lists \texttt{dolbeault-grothendieck-lemma} as a target. As of 2026-10-29 K2 still needs local exactness and the parameter-dependent solver (\href{https://github.com/TauCetiProject/TauCeti/pull/22585}{\#22585}). Layers 1–3 map onto exactly these items; layer 4 is your project's glue.

The proof needs five algebraic identities, five calculus lemmas, one solver with four properties, and a one-line induction step. Indices here are 1-based.

\textbf{Notation.} e\textsubscript{j} is the j-th unit vector of ℂⁿ and dz̄\textsubscript{j}(v) = conj(v\textsubscript{j}). For f with values in a complex Banach space, ∂̄\textsubscript{j} f = ½(D\textsubscript{e\_j} f + i·D\textsubscript{i·e\_j} f), where D\textsubscript{v} is the real directional derivative. The contraction is ι\textsubscript{v} a = a(v, ·). For a real 1-form φ with complex values:

\$\$ (\textbackslash{}varphi\textbackslash{}wedge a)(v\textsubscript{0},\textbackslash{}dots,v\textsubscript{q})=\textbackslash{}sum\textsubscript{i=0}\textsuperscript{q}(-1)\textsuperscript{i}\textbackslash{},\textbackslash{}varphi(v\textsubscript{i})\textbackslash{},a(v\textsubscript{0},\textbackslash{}dots,\textbackslash{}widehat\{v\textsubscript{i}\},\textbackslash{}dots,v\textsubscript{q}) \$\$

The flat Dolbeault operator is

\$\$ \textbackslash{}bar\textbackslash{}partial\textbackslash{}omega=\textbackslash{}sum\textsubscript{j=1}\textsuperscript{n} d\textbackslash{}bar z\textsubscript{j}\textbackslash{}wedge\textbackslash{}bar\textbackslash{}partial\textsubscript{j}\textbackslash{}omega \$\$

\paragraph{A. Pointwise algebra}
\begin{itemize}
\item \textbf{A1 (contraction of a wedge).} ι\textsubscript{v}(φ ∧ a) = φ(v)·a − φ ∧ ι\textsubscript{v} a. Direct from the formula above.
\item \textbf{A2.} ι\textsubscript{v} ι\textsubscript{v} = 0 and ι\textsubscript{v} ι\textsubscript{w} = −ι\textsubscript{w} ι\textsubscript{v}. Immediate from alternation.
\item \textbf{A3.} φ ∧ φ ∧ a = 0 and φ ∧ ψ ∧ a = −ψ ∧ φ ∧ a.
\item \textbf{A4 (type).} If a is of type (0,q), so are ι\textsubscript{v} a and dz̄\textsubscript{j} ∧ a (types (0,q−1) and (0,q+1)). dz̄\textsubscript{j} is conjugate-linear, and the defining relation is linear in each slot.
\item \textbf{A5 (expansion).} If a has type (0,q+1), then a(v, w) = Σ\textsubscript{j} conj(v\textsubscript{j})·(ι\textsubscript{e\_j} a)(w). Write v = Σ\textsubscript{j} (Re v\textsubscript{j})e\textsubscript{j} + (Im v\textsubscript{j})(i e\textsubscript{j}) and use a(i e\textsubscript{j}, ·) = −i·a(e\textsubscript{j}, ·). Corollary: if ι\textsubscript{e\_j} a = 0 for every j, then a = 0.
\end{itemize}
\paragraph{B. Calculus on ℂⁿ}
Let U ⊆ ℂⁿ be open and ω smooth on U.

\begin{itemize}
\item \textbf{B1 (naturality).} For a continuous ℂ-linear L, ∂̄\textsubscript{j}(L ∘ ω) = L ∘ ∂̄\textsubscript{j} ω. In particular ∂̄\textsubscript{j} commutes with ι\textsubscript{v}, with dz̄\textsubscript{i} ∧ (·), and with every map defining type.
\item \textbf{B2.} ∂̄\textsubscript{i} ∂̄\textsubscript{j} ω = ∂̄\textsubscript{j} ∂̄\textsubscript{i} ω on U. Use symmetry of the second real derivative.
\item \textbf{B3 (∂̄² = 0).} Expand ∂̄∂̄ω = Σ\textsubscript{i,j} dz̄\textsubscript{i} ∧ dz̄\textsubscript{j} ∧ ∂̄\textsubscript{i}∂̄\textsubscript{j} ω with B1. The coefficient is symmetric in (i,j) by B2 and the prefix antisymmetric by A3, so the sum equals its own negative.
\item \textbf{B4 (contraction of ∂̄).} ι\textsubscript{e\_m} ∂̄ω = ∂̄\textsubscript{m} ω − ∂̄(ι\textsubscript{e\_m} ω). Apply A1 termwise with dz̄\textsubscript{j}(e\textsubscript{m}) = δ\textsubscript{jm}, then B1.
\item \textbf{B5 (closed ⇒ holomorphic in the tail).} Suppose ω involves only dz̄\textsubscript{1}, …, dz̄\textsubscript{k} on U and ∂̄ω = 0 on U. Then ∂̄\textsubscript{m} ω = 0 on U for every m > k. Indeed ι\textsubscript{e\_m} ω ≡ 0 on the open set U, so its ∂̄ vanishes there, and B4 gives ∂̄\textsubscript{m} ω = ι\textsubscript{e\_m} ∂̄ω = 0.
\item \textbf{B6 (slices).} ∂̄\textsubscript{k} F(z) is the ∂̄-derivative at z\textsubscript{k} of w ↦ F(z with z\textsubscript{k} replaced by w). Chain rule along \texttt{update z k}.
\end{itemize}
\paragraph{C. The parametric solver}
Fix x, k and R > 0 with the closed sup-ball (a polydisc) of radius 2R around x inside U. Choose a cutoff χ : ℂ → ℝ that equals 1 on B(x\textsubscript{k}, R) and has support in the closed disc B̄(x\textsubscript{k}, 3R/2). Write z[k ↦ ζ] for z with its k-th coordinate replaced by ζ, and set

\$\$ G(z)=\textbackslash{}frac\{1\}\{\textbackslash{}pi\}\textbackslash{}int\textsubscript{\textbackslash{}mathbb C}\textbackslash{}frac\{1\}\{t\}\textbackslash{},\textbackslash{}chi(z\textsubscript{k}-t)\textbackslash{},f\textbackslash{}bigl(z[k\textbackslash{}mapsto z\textsubscript{k}-t]\textbackslash{}bigr)\textbackslash{},dA(t) \$\$

This is the Mathlib convolution (K ⋆ g\textsubscript{z})(z\textsubscript{k}) with K(t) = t⁻¹ and g\textsubscript{z}(ζ) = χ(ζ)·f(z[k ↦ ζ]). For z in B(x,R) and ζ in the support of χ, the point z[k ↦ ζ] lies in B(x, 2R) ⊆ U. So g is jointly smooth on B(x,R) × ℂ and supported in a fixed compact set in ζ.

\begin{itemize}
\item \textbf{C1 (smoothness).} G is smooth on B(x,R), by \texttt{contDiffOn\_convolution\_right\_with\_param\_comp} with v(z) = z\textsubscript{k}.
\item \textbf{C2 (solves the k-th equation).} ∂̄\textsubscript{k} G = f on B(x,R). By B6, the slice through z is 2T h, where h(ζ) = χ(ζ)f(z[k ↦ ζ]) is smooth with compact support. TauCeti's Pompeiu converse gives ∂̄(2T h) = h, and h(z\textsubscript{k}) = f(z) because χ(z\textsubscript{k}) = 1.
\item \textbf{C3 (keeps holomorphy).} If j ≠ k and ∂̄\textsubscript{j} f = 0 on U, then ∂̄\textsubscript{j} G = 0 on B(x,R). Differentiate under the integral (\texttt{hasFDerivAt\_convolution\_right\_with\_param}); the integrand's ∂̄\textsubscript{j} derivative is χ·(∂̄\textsubscript{j} f)(z[k ↦ ζ]) = 0.
\item \textbf{C4 (naturality).} L(G\textsubscript{f}(z)) = G\textsubscript{L∘f}(z) for continuous ℂ-linear L, by \texttt{integral\_comp\_comm}. Hence if L ∘ f = 0 on U, then L ∘ G = 0 everywhere.
\end{itemize}
\paragraph{D. The induction}
\textbf{Claim P(k).} Let ω be a smooth (0,q+1)-form on an open U that is closed and involves only dz̄\textsubscript{1}, …, dz̄\textsubscript{k} on U. Then for each x ∈ U there are an open U′ with x ∈ U′ ⊆ U and a smooth (0,q)-form η on U′ with ∂̄η = ω on U′.

\textbf{Base, k = 0.} ω has degree q+1 ≥ 1 and all contractions vanish, so ω = 0 by A5. Take η = 0.

\textbf{Step, k ≥ 1.} Put α := ι\textsubscript{e\_k} ω.

\begin{enumerate}
\item α has type (0,q) (A4) and involves only dz̄\textsubscript{1}, …, dz̄\textsubscript{k−1}: ι\textsubscript{e\_m} ι\textsubscript{e\_k} ω = −ι\textsubscript{e\_k} ι\textsubscript{e\_m} ω is 0 for m > k, and is 0 for m = k by A2.
\item ∂̄\textsubscript{m} α = ι\textsubscript{e\_k} ∂̄\textsubscript{m} ω = 0 on U for m > k (B1, B5).
\item Let G be the solver of C applied to α. By C1–C4, G is smooth on B(x,R), has type (0,q), involves only dz̄\textsubscript{1}, …, dz̄\textsubscript{k−1} everywhere, satisfies ∂̄\textsubscript{k} G = α, and satisfies ∂̄\textsubscript{m} G = 0 for m > k.
\item Put ω′ := ω − ∂̄G on B(x,R). It is closed by B3 and has type (0,q+1) by A4 and B1.
\item ω′ involves only dz̄\textsubscript{1}, …, dz̄\textsubscript{k−1}. For m ≥ k, B4 gives ι\textsubscript{e\_m} ω′ = ι\textsubscript{e\_m} ω − ∂̄\textsubscript{m} G + ∂̄(ι\textsubscript{e\_m} G), and ι\textsubscript{e\_m} G ≡ 0. For m = k this is α − α = 0; for m > k it is 0 − 0.
\item P(k−1) applied to ω′ on B(x,R) gives η′ on U′ ∋ x. Then η := η′ + G satisfies ∂̄η = ω on U′.
\end{enumerate}
\textbf{Theorem.} Every (0,q+1)-form involves only dz̄\textsubscript{1}, …, dz̄\textsubscript{n}, so P(n) is the flat Dolbeault–Grothendieck lemma.

\begin{theorem}[Z4.a — Cauchy kernel]
\label{Z4.a}
\lean{Cohesive.locallyIntegrable_inv_complex}
\leanok
The Cauchy kernel t ↦ t\textsuperscript{−1} is locally integrable on ℂ with respect to Lebesgue measure.
\par\smallskip{\footnotesize Source: blueprint.md §11 v1.14.}
\end{theorem}
\begin{proof}
\leanok
See the Lean proof and the cited items.
\paragraph{Blueprint notes.}
(\texttt{Analysis/CauchyKernel.lean}): \texttt{LocallyIntegrable (fun t : ℂ ↦ t⁻¹)}, from Mathlib's \texttt{locallyIntegrable\_of\_norm\_le\_rpow} (α = 1 < 2 = dim\textsubscript{ℝ} ℂ).
\end{proof}

\begin{theorem}[Z4.b — One-variable ∂̄ (Cauchy–Pompeiu)]
\label{Z4.b}
\lean{Cohesive.dbar_cauchyTransform}
\leanok
Let F be a complex Banach space and h : ℂ → F real C¹ with compact support. Let u(w) = π\textsuperscript{−1} ∫\textsubscript{ℂ} t\textsuperscript{−1} h(w − t) dA(t). Then for every w ∈ ℂ, ½ (Du(w)(1) + i·Du(w)(i)) = h(w), i.e. ∂u/∂w̄ = h in the normalization of F3.
\par\smallskip{\footnotesize Source: blueprint.md §11 v1.14.}
\end{theorem}
\begin{proof}
\leanok
\uses{Z4.a}
See the Lean proof and the cited items.
\paragraph{Blueprint notes.}
(\texttt{Analysis/Pompeiu.lean}): for \texttt{h : ℂ → Λ} (Λ a complex Banach space), \texttt{ContDiff ℝ 1 h}, compact support, \texttt{u w = π⁻¹ ∫ t⁻¹ • h (w − t)} satisfies \texttt{antiPart (fderiv ℝ u w) (toConj 1) = h w}. Proof: \texttt{Du = π⁻¹(t⁻¹ ⋆ Dh)} (Mathlib convolution derivative); polar coordinates around w: the angular term integrates to 0 over each circle, the radial term to −h(w) along each ray, so ∫ ζ⁻¹ ∂̄h(w+ζ) dA = −π h(w). Proposed: port TauCeti \texttt{Pompeiu.lean} (Apache 2.0, with attribution in the module docstring) — needs the human's decision.
\end{proof}

\begin{theorem}[Z4.c — Parametric solver]
\label{Z4.c}
\lean{Cohesive.contDiff_cauchySolve, Cohesive.hasFDerivAt_cauchySolve, Cohesive.map_cauchySolve, Cohesive.dbar_cauchySolve}
\leanok
\uses{Z4.b}
Let E be a complex normed space and Λ a complex Banach space. For v ∈ E and F : E → Λ let u(z) = cauchySolve v F (z) = π\textsuperscript{−1} ∫\textsubscript{ℂ} t\textsuperscript{−1} F(z − t v) dA(t). Let v ≠ 0 and F have compact support.

\begin{itemize}
\item If F is real C\textsuperscript{∞}, then u is real C\textsuperscript{∞} on E.
\item If F is real C¹, then Du(z) = cauchySolve v (DF)(z) for every z.
\item If F is continuous and L : Λ → Λ′ is continuous ℂ-linear (Λ′ Banach), then L(u(z)) = cauchySolve v (L ∘ F)(z).
\item If F is real C¹, then ½ (Du(z)(v) + i·Du(z)(i v)) = F(z) for every z (∂̄ along v solves the equation).
\end{itemize}
\par\smallskip{\footnotesize Source: blueprint.md §11 v1.14.}
\end{theorem}
\begin{proof}
\leanok
See the Lean proof and the cited items.
\paragraph{Blueprint notes.}
(\texttt{Analysis/CauchySolve.lean}): \texttt{cauchySolve v F z := (π⁻¹ : ℂ) • ∫ t : ℂ, t⁻¹ • F (z − t • v)} for \texttt{F : E → Λ} C\textsuperscript{∞} with compact support: (i) C\textsuperscript{∞} on E with \texttt{fderiv (cauchySolve v F) = cauchySolve v (fderiv ℝ F)} (parametric convolution lemmas, g(z,ζ) = F(z + (ζ − ℓ z)•v), ℓ v = 1); (ii) \texttt{L (cauchySolve v F z) = cauchySolve v (L ∘ F) z} for ℂ-linear CLMs L; (iii) for v ≠ 0, \texttt{dbarPart (cauchySolve v F) z (toConj v) = F z} (slice through z, then Z4.b).
\end{proof}

\begin{theorem}[Z4.d — Contraction]
\label{Z4.d}
\lean{Cohesive.eq_zero_of_ctr_basis, Cohesive.dbarPart_comp_clm, Cohesive.ctr_dbar, Cohesive.ctr_dbar_zero, Cohesive.dbarPart_eq_zero_of_dbar_eq_zero}
\leanok
Let E, F be complex normed spaces. For v ∈ Ē, the contraction ι\textsubscript{v} : Λ\textsuperscript{0,q+1}(E; F) → Λ\textsuperscript{0,q}(E; F) is (ι\textsubscript{v} a)(w\textsubscript{1}, …, w\textsubscript{q}) = a(v, w\textsubscript{1}, …, w\textsubscript{q}).

\begin{itemize}
\item If ι\textsubscript{b\_i} a = 0 for every vector of a basis b of Ē, then a = 0.
\item D″(L ∘ ω)(z) = L ∘ D″ω(z) for continuous ℂ-linear maps L between fibres, ω real-differentiable at z.
\item For ω of degree q + 1, real-differentiable at z: ι\textsubscript{v} ∂̄ω(z) = D″ω(z)(v) − ∂̄(ι\textsubscript{v} ω)(z). For ω of degree 0: ι\textsubscript{v} ∂̄ω(z) = D″ω(z)(v).
\item Say a involves only dz̄\textsubscript{0}, …, dz̄\textsubscript{k−1} (for a basis b of Ē indexed by \{0, …, n−1\}) if ι\textsubscript{b\_j} a = 0 for all j ≥ k. If U ⊆ E is open, ω (degree q + 1) is real-differentiable on U, involves only dz̄\textsubscript{0}, …, dz̄\textsubscript{k−1} at every point of U and ∂̄ω = 0 on U, then D″ω(z)(b\textsubscript{j}) = 0 for all z ∈ U and j ≥ k.
\end{itemize}
\par\smallskip{\footnotesize Source: blueprint.md §11 v1.14.}
\end{theorem}
\begin{proof}
\leanok
\uses{F3}
See the Lean proof and the cited items.
\paragraph{Blueprint notes.}
(\texttt{Forms/Contraction.lean}): contraction = \texttt{curryLeft}; anticommutation, \texttt{curryLeft a (toConj v) = 0} on a basis ⇒ \texttt{a = 0}; \texttt{dbarPart} commutes with CLMs; B4 \texttt{curryLeft (dbar ω z) v = dbarPart ω z v − dbar (fun y ↦ curryLeft (ω y) v) z} (degree q + 1) and \texttt{curryLeft (dbar ω z) v = dbarPart ω z v} (degree 0); \texttt{InvolvesOnly b k a} (contractions with \texttt{b j}, j ≥ k, vanish); B5.
\end{proof}

\begin{theorem}[Z4.e — Flat lemma]
\label{Z4.e}
\lean{Cohesive.exists_dbar_eq, Cohesive.conjBasis, Cohesive.dbarPart_cauchySolve, Cohesive.dbarPart_cauchySolve_comm, Cohesive.involvesOnly_cauchySolve, Cohesive.exists_cutoff, Cohesive.contDiff_cutoff_smul}
\leanok
Let E be a finite-dimensional complex normed space with a fixed basis e\textsubscript{0}, …, e\textsubscript{n−1} (also a basis of Ē). Flat Dolbeault–Grothendieck lemma: let U ⊆ E be open, x ∈ U, and β : E → Λ\textsuperscript{0,q+1}(E; ℂ) real C\textsuperscript{∞} on U with ∂̄β = 0 on U. Then there are an open U′ with x ∈ U′ ⊆ U and η : E → Λ\textsuperscript{0,q}(E; ℂ), real C\textsuperscript{∞} on all of E, with ∂̄η = β on U′. Ingredients:

\begin{itemize}
\item for v ≠ 0 and F : E → Λ\textsuperscript{0,q}(E; ℂ) real C¹ with compact support, D″(cauchySolve v F)(z)(v) = F(z) and D″(cauchySolve v F)(z)(w) = cauchySolve v (D″F(·)(w))(z);
\item if every value of F involves only dz̄\textsubscript{0}, …, dz̄\textsubscript{k−1}, so does every value of cauchySolve v F;
\item cutoff: for x ∈ U and an index k there are a real C\textsuperscript{∞} ψ with compact support in U, ψ = 1 on an open W ∋ x, and an open slab S ⊇ W, stable under z ↦ z − t e\textsubscript{k} (t ∈ ℂ), on which the derivatives of ψ in the directions e\textsubscript{j}, i e\textsubscript{j} (j ≠ k) vanish; ψα is C\textsuperscript{∞} with compact support for every α that is C\textsuperscript{∞} on U.
\end{itemize}
\par\smallskip{\footnotesize Source: blueprint.md §11 v1.14.}
\end{theorem}
\begin{proof}
\leanok
\uses{F4, Z4.c, Z4.d}
See the Lean proof and the cited items.
\paragraph{Blueprint notes.}
(\texttt{Forms/FlatGrothendieck.lean}): cutoff ψ (S2), induction step (split on q = 0 / q + 1), \texttt{exists\_dbar\_eq} on open subsets of E with η globally C\textsuperscript{∞}.
\end{proof}

\begin{theorem}[Z4.f — Chart forms]
\label{Z4.f}
\lean{Cohesive.chartOpen, Cohesive.chartForm, Cohesive.chartExpr_chartForm, Cohesive.isDolbeault_of_chartExpr, Cohesive.isDolbeault_chartForm, Cohesive.isOfDegree_chartForm}
\leanok
Let E be a finite-dimensional complex normed space and X a complex manifold modelled on E. For x ∈ X, φ = chartAt(x), an open W ⊆ E and η : E → Λ\textsuperscript{0,q}(E; ℂ), the chart form of η is the form on U′ = φ.source ∩ φ\textsuperscript{−1}(W) whose value at y is η(φ(y)), placed in degree q and pulled back along the derivative of φ ∘ chartAt(y)\textsuperscript{−1} at chartAt(y)(y).

\begin{itemize}
\item If W ⊆ φ.target, the expression of the chart form in φ is η on W.
\item (Converse of F9.a.) If U ⊆ e.source for an atlas chart e and the expression of a form α on U in e is real C\textsuperscript{∞} at e(y) for all y ∈ U, then α is a Dolbeault form.
\item The chart form has degree q; if W ⊆ φ.target and η is real C\textsuperscript{∞} on W, it is a Dolbeault form.
\end{itemize}
\par\smallskip{\footnotesize Source: blueprint.md §11 v1.14.}
\end{theorem}
\begin{proof}
\leanok
\uses{F9}
See the Lean proof and the cited items.
\paragraph{Blueprint notes.}
(\texttt{Forms/ChartForm.lean}): from a flat form η on W ⊆ (chartAt x).target, a Dolbeault form on \texttt{(chartAt x).source ∩ (chartAt x)⁻¹' W} whose expression in \texttt{chartAt x} is η; injectivity of \texttt{chartExpr} at points of the chart. Needed private lemmas of \texttt{Forms/SheafBasic.lean}, \texttt{Forms/DolbeaultFree.lean} become public (visibility only).
\end{proof}

\begin{theorem}[Z4.g — H1: Dolbeault–Grothendieck lemma]
\label{Z4.g}
\lean{Cohesive.dolbeaultGrothendieck, Cohesive.flatExpr}
\leanok
Let E be a finite-dimensional complex normed space and X a complex manifold modelled on E (no compactness or Hausdorff assumption). Then X satisfies the Dolbeault–Grothendieck lemma H1. Proof outline: the flat expression of degree r of a form α in chartAt(x) (the degree-r component of its chart expression) is C\textsuperscript{∞} and ∂̄-closed near chartAt(x)(x); the flat lemma (Z4.e) solves ∂̄η = it, and the chart form of η (Z4.f) is the required local primitive.
\par\smallskip{\footnotesize Source: blueprint.md §11 v1.14.}
\end{theorem}
\begin{proof}
\leanok
\uses{H1, Z4.e, Z4.f}
See the Lean proof and the cited items.
\paragraph{Blueprint notes.}
(\texttt{Forms/DolbeaultGrothendieck.lean}): \texttt{theorem dolbeaultGrothendieck [IsManifold 𝓘(ℂ, E) ∞ X] : DolbeaultGrothendieck E X}.
\end{proof}

\begin{theorem}[Z4.h — Corollary]
\label{Z4.h}
\lean{Cohesive.blockPerf_of_localGaugeLemma, main_theorem_of_localGaugeLemma}
\leanok
\uses{H1}
Let E be a finite-dimensional complex normed space and X a compact Hausdorff complex manifold modelled on E, satisfying the local gauge lemma H2. Then J : Ho(P\textsubscript{A}) ⥤ D(O\textsubscript{X}), A = 𝒜\textsubscript{X}(X), is full and faithful with essential image D\textsubscript{perf}(X, O\textsubscript{X}), for the standard derived category and for any other \texttt{HasDerivedCategory} instance. (H1 holds by Z4.g.)
\par\smallskip{\footnotesize Source: blueprint.md §11 v1.14.}
\end{theorem}
\begin{proof}
\leanok
\uses{M, Z4.g}
See the Lean proof and the cited items.
\paragraph{Blueprint notes.}
(\texttt{MainH1.lean}): \texttt{MainStatement} without H1, as a theorem only (no new statement to freeze; a new main statement follows after Z5). \texttt{\#print axioms} checks.
\end{proof}
